\documentclass[preview]{standalone}

\usepackage{amsmath}
\usepackage{amssymb}
\usepackage{bettelini}
\usepackage{stellar}
\usepackage{definitions}

\begin{document}

\id{measuretheory-complete-measures}
\genpage

\section{Complete measures}

\begin{snippet}{complete-measure-motivation}
    Let \((X,\Sigma,\mu)\) be a measure space, let \(N\in\Sigma\) satisfy
    \(\mu(N)=0\), and let \(E\subseteq N\). It may happen that \(E\notin\Sigma\).
    This is inconvenient in integration: changing a function on a subset of a
    null set should not destroy measurability. A complete measure space rules out
    this phenomenon.
\end{snippet}

\begin{snippetdefinition}{complete-measure-definition}{Complete measure}
    A \measurespace \((X,\Sigma,\mu)\) is \emph{complete} if, for every
    \(N\in\Sigma\) with \(\mu(N)=0\), every subset \(E\subseteq N\) belongs to
    \(\Sigma\).
\end{snippetdefinition}

\begin{plainhtml}
    By monotonicity, every measurable subset of a null set also has measure zero.
\end{plainhtml}

\section{Completion}

\begin{snippetdefinition}{measure-completion-definition}{Completion of a measure space}
    Let \((X,\Sigma,\mu)\) be a \measurespace. Its \emph{completion} is the
    family
    \[
        \Sigma^\mu
        =\left\{E\subseteq X\suchthat
        \begin{array}{c}
            \text{there exist }A,B\in\Sigma\text{ such that}\\
            A\subseteq E\subseteq B
            \text{ and }\mu(B\difference A)=0
        \end{array}
        \right\}.
    \]
    For \(E\in\Sigma^\mu\), define
    \[
        \overline\mu(E)=\mu(A)=\mu(B),
    \]
    where \(A\) and \(B\) are any measurable sets satisfying the displayed
    conditions.
\end{snippetdefinition}

\begin{snippetproposition}{measure-completion-proposition}{Properties of the completion}
    Let \((X,\Sigma,\mu)\) be a \measurespace and let
    \((X,\Sigma^\mu,\overline\mu)\) be its completion. Then:
    \begin{enumerate}
        \item \(\overline\mu\) is well defined;
        \item \(\Sigma^\mu\) is a \(\sigma\)-algebra containing \(\Sigma\);
        \item \(\overline\mu\) is a complete measure on \(\Sigma^\mu\) and
        extends \(\mu\);
        \item \(\Sigma^\mu\) is the smallest \(\sigma\)-algebra containing
        \(\Sigma\) on which \(\mu\) admits a complete extension.
    \end{enumerate}
\end{snippetproposition}

\begin{snippetproof}{measure-completion-proposition-proof}{measure-completion-proposition}{Properties of the completion}
    Suppose first that
    \[
        A\subseteq E\subseteq B,
        \qquad
        A_1\subseteq E\subseteq B_1,
    \]
    where all four bounding sets are in \(\Sigma\) and
    \(\mu(B\difference A)=\mu(B_1\difference A_1)=0\). Since
    \[
        A\difference A_1\subseteq B_1\difference A_1,
        \qquad
        A_1\difference A\subseteq B\difference A,
    \]
    monotonicity gives
    \(\mu(A\difference A_1)=\mu(A_1\difference A)=0\). Finite additivity then
    yields
    \[
        \mu(A)=\mu(A\intersection A_1)=\mu(A_1).
    \]
    The same argument gives \(\mu(B)=\mu(A)\), so \(\overline\mu\) is well defined.

    Every \(A\in\Sigma\) belongs to \(\Sigma^\mu\) by choosing both bounding
    sets equal to \(A\). If \(A\subseteq E\subseteq B\) witnesses
    \(E\in\Sigma^\mu\), then
    \[
        X\difference B\subseteq X\difference E\subseteq X\difference A,
        \qquad
        \mu\bigl((X\difference A)\difference(X\difference B)\bigr)
        =\mu(B\difference A)=0.
    \]
    Hence \(X\difference E\in\Sigma^\mu\). If
    \(A_n\subseteq E_n\subseteq B_n\) witnesses \(E_n\in\Sigma^\mu\), set
    \(A=\bigcup_nA_n\) and \(B=\bigcup_nB_n\). Then
    \[
        A\subseteq\bigcup_nE_n\subseteq B,
        \qquad
        B\difference A\subseteq\bigcup_n(B_n\difference A_n).
    \]
    Countable subadditivity gives \(\mu(B\difference A)=0\), so
    \(\Sigma^\mu\) is a \(\sigma\)-algebra.

    To prove countable additivity, let \(E_n\in\Sigma^\mu\) be pairwise disjoint
    and choose \(A_n\subseteq E_n\subseteq B_n\) as above. The measurable sets
    \(A_n\) are pairwise disjoint. With \(A=\bigcup_nA_n\), the set
    \((\bigcup_nE_n)\difference A\) is contained in the measurable null set
    \(\bigcup_n(B_n\difference A_n)\). Consequently,
    \[
        \overline\mu\left(\bigcup_nE_n\right)
        =\mu(A)
        =\sum_n\mu(A_n)
        =\sum_n\overline\mu(E_n).
    \]
    Thus \(\overline\mu\) is a measure, and it agrees with \(\mu\) on \(\Sigma\).

    If \(E\in\Sigma^\mu\) and \(\overline\mu(E)=0\), choose
    \(A\subseteq E\subseteq B\) as above. Then \(\mu(B)=0\), and for every
    \(N\subseteq E\) the inclusions \(\emptyset\subseteq N\subseteq B\) show
    that \(N\in\Sigma^\mu\) and \(\overline\mu(N)=0\). Hence the completion is
    complete.

    Finally, let \(\mathcal T\) be a \(\sigma\)-algebra containing \(\Sigma\)
    on which \(\mu\) has a complete extension \(\nu\). If
    \(A\subseteq E\subseteq B\) and \(\mu(B\difference A)=0\), then
    \(E=A\union(E\difference A)\), where
    \(E\difference A\subseteq B\difference A\). Completeness gives
    \(E\difference A\in\mathcal T\), so \(E\in\mathcal T\). Therefore
    \(\Sigma^\mu\subseteq\mathcal T\), proving minimality.
\end{snippetproof}

\begin{snippet}{borel-lebesgue-completion-preview}
    The restriction of Lebesgue measure to the Borel \(\sigma\)-algebra is not
    complete. Its completion is Lebesgue measure on the Lebesgue
    \(\sigma\)-algebra, and therefore
    \[
        \mathcal B(\realnumbers)\subseteq\mathcal L(\realnumbers).
    \]
\end{snippet}

\end{document}
